\chapter{Gõ Nhẹ Mà Trúng}
\chapterintro{Sửa mà không làm lớp cụt hứng}{Một câu góp ý khéo có thể giữ được cả không khí lẫn kỷ luật.}

\section{Em thông minh lắm, giờ dùng vào bài giùm thầy cô}
\begin{manthoai}{Màn thoại}
\textbf{Thầy/Cô:} Em thông minh lắm, giờ dùng vào bài giùm thầy cô.\\
\textbf{Học sinh:} Dạ.
\end{manthoai}
\begin{dungkhinào}
Hợp với học sinh lanh, nhiều năng lượng, hay pha trò lệch nhịp.
\end{dungkhinào}
\begin{nhipthem}
Sau câu này nên giao việc ngay: trả lời, viết bảng, tóm ý, phát giấy.
\end{nhipthem}
\ngancach

\section{Nói nhỏ thôi, ý em không nhỏ đâu}
\begin{manthoai}{Màn thoại}
\textbf{Thầy/Cô:} Nói nhỏ thôi, ý em không nhỏ đâu.\\
\textbf{Học sinh:} Dạ.
\end{manthoai}
\begin{dungkhinào}
Dùng khi học sinh quá hào hứng, nói chen hoặc nói to.
\end{dungkhinào}
\begin{nhipthem}
Câu này giữ được sự trân trọng với nội dung học sinh nói ra, nên các em dễ nghe hơn hẳn kiểu \textit{``Im lặng''} khô cứng.
\end{nhipthem}
\ngancach

\section{Em đang vui, nhưng cho bài học vui ké với}
\begin{manthoai}{Màn thoại}
\textbf{Thầy/Cô:} Em đang vui, nhưng cho bài học vui ké với.\\
\textbf{Học sinh:} Dạ, em quay lại ạ.
\end{manthoai}
\begin{dungkhinào}
Hợp cho lớp vui quá đà mà thầy cô chưa muốn cắt mạnh tay.
\end{dungkhinào}
\begin{nhipthem}
Nếu dùng xong mà lớp vẫn trôi, phải chuyển ngay sang hoạt động cụ thể chứ không dừng ở câu nói dí dỏm.
\end{nhipthem}
\ngancach

\section{Em nhanh quá rồi, chờ lớp một nhịp giùm thầy cô}
\begin{manthoai}{Màn thoại}
\textbf{Thầy/Cô:} Em nhanh quá rồi, chờ lớp một nhịp giùm thầy cô.\\
\textbf{Học sinh:} Dạ.
\end{manthoai}
\begin{dungkhinào}
Hợp với học sinh lanh, nói chen hoặc lao lên quá sớm khiến bạn khác khó theo.
\end{dungkhinào}
\begin{nhipthem}
Câu này không dập sự nhiệt tình, chỉ chỉnh tốc độ để lớp vẫn đi cùng nhau.
\end{nhipthem}
\ngancach

\section{Câu đùa này giữ lại cho giờ ra chơi nhé}
\begin{manthoai}{Màn thoại}
\textbf{Thầy/Cô:} Câu đùa này giữ lại cho giờ ra chơi nhé, còn giờ mình giữ ý chính trước.\\
\textbf{Học sinh:} Dạ.
\end{manthoai}
\begin{dungkhinào}
Dùng khi lớp bắt đầu vui quá nhịp nhưng thầy cô vẫn muốn giữ bầu không khí dễ chịu.
\end{dungkhinào}
\begin{nhipthem}
Rất hiệu quả nếu sau câu này thầy cô giao ngay một việc cụ thể để chuyển năng lượng về bài.
\end{nhipthem}
